\section{Optional extension: time-adaptive smoothness}
\label{sec:optional_dynamic}

The pooled criterion in \cref{eq:F_S} is complete without any
time-adaptive mechanism. If an analyst instead assumes that recent
changes in the desirable smoothing regime are informative about
future forecast performance, the local minima of successive
chronological forecast-loss surfaces can be tracked through time.
The resulting branch state is an optional hyperparameter-selection
layer, not a new trend estimator or a requirement for pooled
forecast cross-validation.

\subsection{Tracking local minima through time}

To retain that organization, we connect nearby local minima across adjacent forecast origins. A current minimum $S_{j,m}$ may continue a previous branch when
\begin{equation}
\label{eq:epsilon_tracking}
|S_{j,m}-S_{j,m-1}|
\le
\varepsilon,
\end{equation}
subject to one-to-one matching at each origin.

The radius $\varepsilon$ does not define a fixed exogenous region such as ``smoothness near 0.9.'' Instead, branch locations are determined by the data through the recovered local minima; $\varepsilon$ only controls whether minima on neighboring surfaces are close enough to be treated as the same temporal branch.

For a tracked branch $j$, let $\ell^{(1)}_{j,m}$ denote its Validation-1 forecast loss at the origin where the local minimum is identified. The trend is then refit through Validation 1 using the same $(d,L,S_{j,m})$, and the immediately following block is scored to obtain $\ell^{(2)}_{j,m}$, the Validation-2 forecast loss.

We collect the branch history as
\begin{equation}
\label{eq:branch_matrix}
V_j
=
\begin{pmatrix}
S_{j,t_1} & \ell^{(1)}_{j,t_1} & \ell^{(2)}_{j,t_1}\\
S_{j,t_2} & \ell^{(1)}_{j,t_2} & \ell^{(2)}_{j,t_2}\\
\vdots & \vdots & \vdots
\end{pmatrix}.
\end{equation}

\subsection{Branch selection and current smoothness}

For this optional extension, the forecasting decision is separated into two steps. First, a historical branch-selection rule
\begin{equation}
\label{eq:branch_selector}
\widehat j_T
=
\psi(V_1,\ldots,V_J)
\end{equation}
chooses among persistent branches using information available before the outer test. A simple baseline prefers branches with maximal support and then minimizes mean historical Validation-2 loss.

Second, a branch-to-smoothness rule
\begin{equation}
\label{eq:phi_rule}
\widehat S_T
=
\phi(V_{\widehat j_T})
\end{equation}
converts the selected branch history into the smoothness used at the current forecast origin.

The simplest rule is the newest tracked minimum,
\begin{equation}
\label{eq:rule_latest}
\phi_{\mathrm{last}}(V_j)=S_{j,T}.
\end{equation}
Alternative stabilizing rules use the recent mean or median,
\begin{equation}
\label{eq:rule_recent_mean}
\phi_{\mathrm{mean},K}(V_j)
=
\frac1K
\sum_{r=0}^{K-1}
S_{j,T-r},
\end{equation}
or
\begin{equation}
\label{eq:rule_recent_median}
\phi_{\mathrm{med},K}(V_j)
=
\operatorname{median}
\{S_{j,T-K+1},\ldots,S_{j,T}\}.
\end{equation}

An alternative weighted mean uses recency alone. Let $r=0$ denote the newest tracked smoothness value, with older values indexed by increasing $r$. Then
\begin{equation}
\label{eq:rule_recency}
\phi_{\mathrm{recency}}(V_j)
=
\frac{
\sum_{r=0}^{R}\rho^r S_{j,T-r}
}{
\sum_{r=0}^{R}\rho^r
},
\qquad
0<\rho<1.
\end{equation}
The newest value receives weight one and older values are exponentially discounted. For interpretability we parameterize $\rho$ by a half-life $H$,
\begin{equation}
\label{eq:recency_half_life}
\rho=2^{-1/H},
\end{equation}
so the weight assigned to a value $H$ tracked origins old is one half of the newest weight.

A loss-weighted rule gives more influence to historically predictive values:
\begin{equation}
\label{eq:rule_validation2}
\phi_{\mathrm{V2}}(V_j)
=
\frac{
\sum_t w_{j,t}S_{j,t}
}{
\sum_t w_{j,t}
},
\qquad
w_{j,t}
=
\frac{1}{\ell^{(2)}_{j,t}+\delta},
\end{equation}
with $\delta>0$ used only for numerical stabilization. Recency may be incorporated through
\begin{equation}
\label{eq:recency_validation2_weights}
w_{j,t}
=
\frac{\rho^{T-t}}{\ell^{(2)}_{j,t}+\delta},
\qquad
0<\rho\le1.
\end{equation}

A distinct class of maps treats the branch smoothness values as a series to be extrapolated. For a branch with values $S_{j,t_1},\ldots,S_{j,t_m}$, fit the auxiliary trend
\begin{equation}
\label{eq:branch_linear_fit}
(\widehat a,\widehat b)
\in
\arg\min_{a,b}
\sum_{u=1}^{m}
\omega_u\bigl(S_{j,t_u}-a-bu\bigr)^2.
\end{equation}
The resulting branch-based smoothness decision is
\begin{equation}
\label{eq:rule_linear_trend}
\phi_{\mathrm{trend}}(V_j)
=
\Pi_{[0,1]}\bigl(\widehat a+\widehat b(m+1)\bigr),
\end{equation}
where $\Pi_{[0,1]}$ truncates to the admissible normalized-smoothness interval. Uniform weights, recent-window weights, or exponentially increasing recency weights lead to different examples. Alternatively, one can extrapolate the branch increments:
\begin{equation}
\label{eq:rule_increment_extrapolation}
\phi_{\Delta,\rho}(V_j)
=
\Pi_{[0,1]}
\left[
S_{j,t_m}
+
\frac{
\sum_{r=0}^{m-2}
\rho^r\bigl(S_{j,t_{m-r}}-S_{j,t_{m-r-1}}\bigr)
}{
\sum_{r=0}^{m-2}\rho^r
}
\right].
\end{equation}
These auxiliary models forecast the next \emph{smoothness value}, not the underlying observations or an old fitted trend. The final observation forecast still follows the fresh-refit procedure below.

The examples above are not intended to define a tournament for one universally best smoothness rule. Their role is to show that the tracked branch state $V_j$ supports a family of admissible decision maps
\begin{equation}
\label{eq:branch_decision_map}
\phi:V_j\mapsto \widehat S_T.
\end{equation}
Different applications may emphasize persistence, recency, historical forecast performance, robustness, or forward extrapolation of the smoothness trajectory. Empirical comparisons describe the behavior of these choices under controlled conditions; they do not make one functional part of the definition of the framework.


The mappings $\psi$ and $\phi$ express different modeling assumptions
about persistence, recency, and changing noise or trend dynamics.
They must be calibrated on historical development data and then
evaluated without retuning on the final outer test. If such
assumptions are unjustified, pooled forecast-CV remains the
self-contained method.
